\documentclass[11pt,a4paper]{article}
\usepackage[left=1in,right=1in,top=0.85in,bottom=0.85in]{geometry}
\usepackage{amsmath,amssymb,amsthm}
\usepackage{microtype}
\usepackage[hidelinks]{hyperref}
\newtheorem{proposition}{Proposition}
\title{Nonclosedness of an ordinary Sobolev restriction
of a magnetic Airy operator}
\author{Kaiwen Jin\thanks{Email: \href{mailto:yui.ui.siki@gmail.com}{\texttt{yui.ui.siki@gmail.com}}.}}
\date{}
\begin{document}
\maketitle
\begin{abstract}
An ordinary Sobolev domain need not give a closed realization of a magnetic
differential expression on an unbounded set. For a half-plane magnetic
complex Airy expression, we construct an explicit graph limit showing that
the restriction printed with the ordinary $H^1_0$ domain is nonclosed for
every fixed real parameter. Under the bounded-inverse definition of the
resolvent, its spectrum is all of $\mathbb C$; if spectrum is reserved for
closed operators, the domain fails the prerequisite. This domain correction
leaves the all-parameter spectral question for the natural closed magnetic
Dirichlet realization unanswered.
\end{abstract}

Almog's half-plane magnetic Airy problem \cite{Almog2015} is printed with an
ordinary $H^1_0$ domain. The spectral problem in Almog--Helffer
\cite[pp.~1--2 and~6]{AlmogHelffer} concerns a closed magnetic realization, whose
Sobolev norm uses covariant derivatives. This note distinguishes these domain
prescriptions by an explicit graph limit for the ordinary restriction.

The construction places increasingly oscillatory functions at large $x$.
Their $y$ frequency cancels the vector potential near the center, making
the graph norm grow more slowly than the ordinary $y$ derivative.
A disjoint sum then converges in graph norm while leaving ordinary $H^1$.

\section{The ordinary Sobolev restriction}
\label{sec:operator}

Let $\Omega=\{(x,y)\in\mathbb R^2:y>0\}$ and, for a fixed $c\in\mathbb R$, write
\begin{equation}
 \mathcal A_c=-\partial_x^2-\left(\partial_y-\frac{ix^2}{2}\right)^2+icy.
 \label{eq:expression}
\end{equation}
All functions are complex-valued. We distinguish this differential expression
from the operator $T_c$ on $L^2(\Omega)$ with domain
\begin{equation}
 D(T_c)=H^1_0(\Omega)\cap
 \{u\in L^2(\Omega):\mathcal A_cu\in L^2(\Omega)\},
 \qquad T_cu=\mathcal A_cu.
 \label{eq:domain}
\end{equation}
The expression in the domain is understood distributionally. Here $H^1_0$
is the closure of $C_c^\infty(\Omega)$ in the \emph{ordinary} Sobolev norm,
which includes both $\|\partial_xu\|_2$ and $\|\partial_yu\|_2$.

For a possibly nonclosed operator $T$ on a Hilbert space $H$, let
$\rho(T)$ consist of the numbers $z\in\mathbb C$ for which
$T-z:D(T)\to H$ is bijective and its inverse is bounded as a map from $H$
to $H$. We then set
\begin{equation}
 \sigma(T)=\mathbb C\setminus\rho(T).
 \label{eq:resolvent}
\end{equation}
If spectrum is defined only for closed operators, the nonclosedness assertion
below says instead that \eqref{eq:domain} does not meet that definition's
prerequisite.

\begin{proposition}
\label{prop:main}
For every $c\in\mathbb R$, the operator $T_c$ in \eqref{eq:domain} is densely
defined and nonclosed. With convention \eqref{eq:resolvent},
$\rho(T_c)=\varnothing$ and $\sigma(T_c)=\mathbb C$.
\end{proposition}

Density follows from $C_c^\infty(\Omega)\subset D(T_c)$ and the density of test
functions in $L^2$ \cite[Corollary~4.23]{Brezis}. Every polynomial coefficient
is bounded on the support of a test function, so the inclusion uses exactly
the domain in \eqref{eq:domain}. The proof of nonclosedness is given next;
the resolvent consequence is established in Section~\ref{sec:consequences}.

\section{An explicit graph limit}
\label{sec:construction}

Choose real functions
$\phi\in C_c^\infty((-1/4,1/4))$ and
$\psi\in C_c^\infty((1,2))$ with $\|\phi\|_2=\|\psi\|_2=1$.
Norms of one-dimensional functions are taken on $\mathbb R$; all other
norms in this section are on $\Omega$.
For $R\geq1$ set
\begin{equation}
 t=\sqrt R(x-R),\qquad k_R=\frac{R^2}{2},\qquad
 v_R(x,y)=R^{1/4}\phi(t)\psi(y)e^{ik_Ry}.
 \label{eq:packet}
\end{equation}
These are smooth compactly supported functions in $\Omega$, and
$\|v_R\|_2=1$ by the substitution $dx=R^{-1/2}dt$.
The carrier frequency $k_R$ cancels the vector potential at $x=R$.
More precisely, on the support of $v_R$,
\begin{equation}
 q_R(x):=k_R-\frac{x^2}{2}
       =-\sqrt R\,t-\frac{t^2}{2R},\qquad
 |q_R(x)|\leq Q\sqrt R,\qquad Q=\frac9{32}.
 \label{eq:mismatch}
\end{equation}
Since $q_R$ is independent of $y$, direct differentiation gives
\begin{equation}
\begin{split}
 \mathcal A_cv_R=R^{1/4}e^{ik_Ry}\bigl[
 &-R\phi''(t)\psi(y)-\phi(t)\psi''(y)
       -2iq_R(x)\phi(t)\psi'(y)\\
 &+q_R(x)^2\phi(t)\psi(y)+icy\phi(t)\psi(y)\bigr].
\end{split}
\label{eq:image}
\end{equation}
Thus
\begin{equation}
 \|T_cv_R\|_2
 \leq R\|\phi''\|_2+\|\psi''\|_2
       +2Q\sqrt R\|\psi'\|_2+Q^2R+2|c|
 \leq C_cR,
 \label{eq:graphbound}
\end{equation}
where one may take
$C_c=\|\phi''\|_2+\|\psi''\|_2+2Q\|\psi'\|_2+Q^2+2|c|$.
It is finite for every fixed $c$ and independent of $R$.
In contrast, the ordinary derivative has norm
\begin{equation}
 \|\partial_yv_R\|_2^2
 =\|\psi'+ik_R\psi\|_2^2
 =\|\psi'\|_2^2+\frac{R^4}{4}.
 \label{eq:ordinary}
\end{equation}
The last equality uses the reality of $\psi$ and $\psi'$, which makes the
cross term zero.

Now let $R_n=4^n$ and $a_n=R_n^{-3/2}$ for $n\geq1$, and put
\begin{equation}
 u_N=\sum_{n=1}^N a_nv_{R_n},\qquad
 u=\sum_{n=1}^\infty a_nv_{R_n},\qquad
 f_c=\sum_{n=1}^\infty a_nT_cv_{R_n}.
 \label{eq:series}
\end{equation}
The $x$ support of $v_{R_n}$ lies in
$I_n=(R_n-1/(4\sqrt{R_n}),R_n+1/(4\sqrt{R_n}))$.
These intervals are pairwise disjoint: successive centers are separated by
$3R_n$, whereas the sum of the two half-widths is $3/(8\sqrt{R_n})$.
Derivatives and multiplication by the coefficients of $\mathcal A_c$ do not
enlarge support, so the images in \eqref{eq:series} are disjoint as well.

Every compact subset of $\Omega$ meets only finitely many of these supports.
The two infinite sums therefore define actual locally smooth functions.
Disjointness and \eqref{eq:graphbound} show that they belong to $L^2(\Omega)$,
with the explicit tails
\begin{align}
 \|u-u_N\|_2^2
 &=\sum_{n>N}4^{-3n}
   =\frac{4^{-3(N+1)}}{1-4^{-3}}\longrightarrow0,
 \label{eq:utail}\\
 \|f_c-T_cu_N\|_2^2
 &\leq C_c^2\sum_{n>N}4^{-n}
   =\frac{C_c^2\,4^{-(N+1)}}{1-4^{-1}}\longrightarrow0.
 \label{eq:ftail}
\end{align}
The squared-norm identities for the infinite sums follow from monotone
convergence applied to the nonnegative disjoint terms.
Local finiteness also gives $\mathcal A_cu=f_c$ distributionally: a compactly
supported test function sees only finitely many terms of \eqref{eq:series}.

Nevertheless $u\notin H^1(\Omega)$. Indeed, if $g\in L^2(\Omega)$ were its
ordinary weak $y$ derivative, then on the open strip $I_n\times(0,\infty)$
the weak derivative identity would give
$g=a_n\partial_yv_{R_n}$ almost everywhere. Uniqueness here is the usual
uniqueness of a locally integrable distribution
\cite[Corollary~4.24]{Brezis}.
For every $M\geq1$, \eqref{eq:ordinary} would then imply
\begin{equation}
 \|g\|_2^2\geq
 \sum_{n=1}^M a_n^2\|\partial_yv_{R_n}\|_2^2
 \geq\frac14\sum_{n=1}^M R_n
 \geq\frac{4^M}{4},
 \label{eq:divergence}
\end{equation}
a contradiction. Thus $u\notin D(T_c)$, although
$u_N\to u$ and $T_cu_N\to f_c$ strongly in $L^2$.
Since every $u_N$ belongs to $C_c^\infty(\Omega)\subset D(T_c)$, this proves
that the graph is not closed. The functions are supported in $1<y<2$; the
failure occurs at spatial infinity, away from the boundary.

\section{The resolvent and the closed magnetic realization}
\label{sec:consequences}

For completeness, a bounded everywhere-defined inverse of a bijective shift
forces an operator to be closed. To see this, suppose $B=(T-z)^{-1}:H\to H$
has these properties, and take $w_j\in D(T)$ with
$w_j\to w$ and $Tw_j\to h$ in $H$. Boundedness gives
\[
 w_j=B(Tw_j-zw_j)\longrightarrow B(h-zw).
\]
Hence $w=B(h-zw)\in D(T)$ and $Tw=h$.
Applied to the nonclosed $T_c$, this proves $\rho(T_c)=\varnothing$ under
\eqref{eq:resolvent}, completing the proof of Proposition~\ref{prop:main}.
The statement $\sigma(T_c)=\mathbb C$ does not assert that every complex
number is an eigenvalue.

The domain distinction is important for interpreting this conclusion.
The closed realization discussed in \cite[pp.~6 and~12,
equations~(2.3), (2.4), and~(2.24)]{AlmogHelffer} is defined using the
magnetic space $H^1_{0,A}(\Omega)$, the closure of $C_c^\infty(\Omega)$ for
the norm
\[
 \left(\|u\|_2^2+\|\partial_xu\|_2^2+
 \left\|\left(\partial_y-\frac{ix^2}{2}\right)u\right\|_2^2\right)^{1/2}.
\]
For $c\ne0$, its form domain also requires $y^{1/2}u\in L^2(\Omega)$;
the operator domain is then obtained by requiring the distributional
expression to lie in $L^2$. This prescription uses the covariant derivative
rather than the ordinary $\partial_y$ norm in \eqref{eq:ordinary}.

The result above concerns only the ordinary restriction \eqref{eq:domain},
which is nonclosed. It identifies a defect in that printed domain, or gives
the convention-dependent spectrum stated in Proposition~\ref{prop:main}.
It provides no conclusion about spectral nonemptiness, eigenvalues, or a
spectral threshold for the distinct natural closed magnetic
Dirichlet realization. That all-parameter spectral question is left
unanswered by this note.

\paragraph{AI assistance.}
OpenAI Codex assisted with the mathematical construction, source checks,
writing and two draft reviews. These AI reviews do not constitute independent
human peer review.

\begin{thebibliography}{9}
\bibitem{Almog2015}
Y.~Almog, \emph{Magnetic Schr\"odinger operator}, AIM problem sheet,
2015, p.~1.
\href{https://nsa.fjfi.cvut.cz/problems/04_AIM_2015/2015-02/2015-02_AIM_Almog_2.pdf}{AIM problem sheet}.

\bibitem{AlmogHelffer}
Y.~Almog and B.~Helffer, \emph{On the spectrum of non-selfadjoint
Schr\"odinger operators with compact resolvent},
\href{https://arxiv.org/abs/1410.5122v1}{arXiv:1410.5122v1}, 2014.

\bibitem{Brezis}
H.~Brezis, \emph{Functional Analysis, Sobolev Spaces and Partial Differential
Equations}, Universitext, Springer, 2011.
\href{https://doi.org/10.1007/978-0-387-70914-7}{doi:10.1007/978-0-387-70914-7}.
\end{thebibliography}
\end{document}
